\input{TAC-801-SHARED-PREAMBLE.tex}
\begin{document}
\nolinenumbers
\begin{singlespace}
Franciscus Dylan Rosario, \\
7624 Melody \\
Rohnert Park, CA 94928 \\
E: soltrinox@gmail.com \\
Plaintiff, In Pro Per \\
\end{singlespace}
\vspace*{9mm}
\begin{tightcenter}
\textbf{SUPERIOR COURT OF THE STATE OF CALIFORNIA} \\
\textbf{COUNTY OF SAN FRANCISCO}
\end{tightcenter}
\vspace*{2.25mm}
\nolinenumbers
\begin{minipage}[t]{3in}
    \raggedright
    \textbf{Franciscus Dylan Rosario,} \\ \textbf{PLAINTIFF}, \\
    \hspace{1cm} v. \\
    \small
    CSAA Insurance Exchange; \\
    Carbone, Smith \& Koyama, LLP; \\
    Michael R. Chambers; \\
    Phillips, Spallas \& Angstadt LLP; \\
    Priya D. Navaratnasingham; \\
    Alberto Reyna; \\
    and DOES 1 through 20, inclusive; \\
    \normalsize
    \textbf{DEFENDANTS}
    \makebox[3in]{\hrulefill}
\end{minipage}%
\hspace{3mm}
\begin{minipage}[t]{0.5pt}
    \vspace{0pt}
    \rule{0.5pt}{3.1in}
\end{minipage}
\hspace{3mm}
\begin{minipage}[t]{3in}
    \raggedright
    \noindent \textbf{Case No.: CGC-25-631801} \\
    \textbf{(Related: CGC-21-594102)} \\
    ~\\
    \textbf{MEMORANDUM OF POINTS AND AUTHORITIES} \\
    \textbf{IN SUPPORT OF LEAVE TO FILE THIRD AMENDED COMPLAINT}
\end{minipage}
\vspace*{6mm}
\linenumbers
\begin{center}\textbf{TABLE OF CONTENTS}\end{center}
\vspace{2mm}
\CourtTOCRow{I.}{\tacctoclink{i-introduction}{I. INTRODUCTION}}{\pageref{toc:i-introduction}}
\CourtTOCRow{II.}{\tacctoclink{ii-statement-of-facts}{II. STATEMENT OF FACTS}}{\pageref{toc:ii-statement-of-facts}}
\CourtTOCRow{III.}{\tacctoclink{iii-legal-standard}{III. LEGAL STANDARD}}{\pageref{toc:iii-legal-standard}}
\CourtTOCRow{IV.}{\tacctoclink{iv-the-tac-cures-all-eight-defense-objections}{IV. THE TAC CURES ALL EIGHT DEFENSE OBJECTIONS}}{\pageref{toc:iv-the-tac-cures-all-eight-defense-objections}}
\CourtTOCRow{V.}{\tacctoclink{v-specific-answers-to-defenses-legal-arguments}{V. SPECIFIC ANSWERS TO DEFENSE'S LEGAL ARGUMENTS}}{\pageref{toc:v-specific-answers-to-defenses-legal-arguments}}
\CourtTOCRow{VI.}{\tacctoclink{vi-no-prejudice-to-defendants}{VI. NO PREJUDICE TO DEFENDANTS}}{\pageref{toc:vi-no-prejudice-to-defendants}}
\CourtTOCRow{VII.}{\tacctoclink{vii-conclusion}{VII. CONCLUSION}}{\pageref{toc:vii-conclusion}}
\vspace{4mm}

\begin{center}\textbf{TABLE OF AUTHORITIES}\end{center}
\vspace{2mm}
\subsubsection*{Cases}
\caseentry{kulchar}{Kulchar v. Kulchar}{(1969) 1 Cal.3d 467}
\caseentry{modnick}{In re Marriage of Modnick}{(1983) 33 Cal.3d 897}
\caseentry{sanders}{Estate of Sanders}{(1985) 40 Cal.3d 607}
\caseentry{westphal}{Westphal v. Westphal}{(1942) 20 Cal.2d 393}
\caseentry{kougasian}{Kougasian v. TMSL, Inc.}{(2004) 118 Cal.App.4th 1028}
\caseentry{stevenot}{In re Marriage of Stevenot}{(1984) 154 Cal.App.3d 1051}
\caseentry{olivera}{Olivera v. Grace}{(1942) 19 Cal.2d 570}
\caseentry{flatley}{Flatley v. Mauro}{(2006) 39 Cal.4th 299}
\caseentry{actionapt}{Action Apartment Ass'n v. City of Santa Monica}{(2007) 41 Cal.4th 1232}
\caseentry{howard_sd}{Howard v. County of San Diego}{(2010) 184 Cal.App.4th 1422}
\caseentry{morgan}{Morgan v. Superior Court}{(1959) 172 Cal.App.2d 527}
\subsubsection*{Statutes}
\statuteentry{ccp473}{Code Civ. Proc.,}{\textsection~473}
\statuteentry{civ2338}{Civ. Code,}{\textsection~2338}
\statuteentry{bpc6128}{Bus. \& Prof. Code,}{\textsection~6128}
\vspace{4mm}
\phantomsection\label{toc:i-introduction}
\section{I. INTRODUCTION}\label{i.-introduction}

Plaintiff FRANCISCUS DYLAN ROSARIO respectfully submits this Memorandum in support of his Motion for Leave to File a Third Amended Complaint (``TAC'').

\textbf{Two-Layer Cure Strategy:} This TAC motion is the \textbf{second layer} of Plaintiff's cure strategy. The \textbf{first layer} is the pending Motion for Leave to File Second Amended Complaint (filed April 21, 2026; hearing May 6, 2026). If the SAC is granted and defense brings Round D challenges, this TAC stands ready as a further cure. The proposed TAC cures the eight pleading objections raised in Defendants' Motion to Strike filed April 22, 2026, and installs a California Supreme Court doctrinal spine that forecloses each defense objection as a matter of binding state Supreme Court law.

\subsection{I.A --- The Controlling Rule}\label{i.a-the-controlling-rule}

The controlling rule is stated by the California Supreme Court in \emph{Westphal v. Westphal} (1942) 20 Cal.2d 393, 397: ``Fraud is extrinsic \ldots{} where the unsuccessful party has been prevented from exhibiting fully his case, by fraud or deception practiced on him by his opponent \ldots{} or where the defendant never had knowledge of the suit, being kept in ignorance by the acts of the plaintiff.'' The Supreme Court has restated that rule three times in terms directly applicable here:

\begin{itemize}
\tightlist
\item
  \emph{Kulchar v. Kulchar} (1969) 1 Cal.3d 467, 471: the ``primary right'' protected is ``a fair adversary hearing.''
\item
  \emph{In re Marriage of Modnick} (1983) 33 Cal.3d 897, 905: extrinsic fraud is a ``broad concept'' reaching ``any extrinsic circumstance'' that has ``prevented a party from fully participating'' in the proceeding.
\item
  \emph{Estate of Sanders} (1985) 40 Cal.3d 607, 614: extrinsic fraud exists where the complained-of conduct ``effectively prevented \ldots{} a fair adversary hearing'' (quoting \emph{Larrabee v. Tracy} (1943) 21 Cal.2d 645, 650--651).
\end{itemize}

Under each of these Supreme Court holdings, the injury is \emph{deprivation of the opportunity to present the case}, not \emph{non-admission of a particular exhibit}. That distinction is dispositive of the defense's anticipated ``it got in on Day 7'' argument (\textsection{} I.F infra).

\subsection{I.B --- The ``Never Fairly Before the Court'' Standard: Denial of the Existence of the Evidentiary Record}\label{i.b-the-never-fairly-before-the-court-standard-denial-of-the-existence-of-the-evidentiary-record}

\emph{In re Marriage of Modnick} (1983) 33 Cal.3d 897, 907, supplies the operative sub-rule. The Supreme Court held that a party who conceals material matter from the tribunal effects a ``denial of the very existence'' of the concealed matter, with the doctrinal consequence that the subject is ``never fairly before the court for adjudication.'' (\emph{Id.} at p.~907, quoting \emph{Clark v. Clark} (1961) 195 Cal.App.2d 373, 381; accord \emph{Jorgensen v. Jorgensen} (1948) 32 Cal.2d 13, 21.) \emph{Modnick} locates the wrong not in the falsity of any particular representation (which would be intrinsic), but in the structural \emph{suppression of subject matter} from the adversary process itself. The TAC pleads that wrong in four linked subparagraphs (TAC ¶¶ 3.5--3.5.4) that mirror the \emph{Modnick} syllogism:

\textbf{1. The Suppressed Subject Matter (TAC ¶ 3.5.1).} The 911 audio of August 18, 2022, the SFPD body-worn camera footage, the CAD entries, and the supplemental police report are not peripheral impeachment. They \emph{constitute} the primary, contemporaneous, government-generated record of the incident on which the defense, the defense declarations, and the Court's structural rulings all depended.

\textbf{2. The Denial (TAC ¶ 3.5.2).} Defendants did not dispute weight; they denied existence. MIL No.~17 (Jan.~17, 2025) represented ``no idea where this 911 call log or audio recording came from.'' The Feb.~28, 2025 colloquy (Hr'g Tr. pp.~291:26--292:7) is the on-record predicate that the trial court extended institutional reliance to defense counsel's representation of non-receipt. Under \emph{Modnick}, a denial of existence is categorically distinct from a credibility dispute over the contents of a known exhibit.

\textbf{3. The Doctrinal Consequence (TAC ¶ 3.5.3).} Because the concealed subject matter was ``never fairly before the court for adjudication'' (\emph{Modnick}, supra, 33 Cal.3d at p.~907), every ruling predicated on the assumed completeness of the August 18, 2022 record---the seven consequential rulings catalogued at \textsection{} I.D infra---was procured in a proceeding from which the decisive evidentiary subject matter had been structurally withheld. When defense counsel's representations were later exposed as false by the April 16, 2025 admission (RT 7, p.~213, ll. 13--15 --- ``They were provided by Dolan''), the \emph{Modnick} defect crystallized: the record of possession and provenance had never been ``fairly before the court'' at the time of any of the seven rulings.

\textbf{4. The \emph{Pico} Line Does Not Reach This Pleading (TAC ¶ 3.5.4).} The rule of \emph{Pico v. Cohn} (1891) 91 Cal. 129, 133--134---that perjury alone is intrinsic---is not implicated. Plaintiff does not allege that Defendants testified falsely about the \emph{contents} of the concealed materials. Plaintiff alleges that Defendants denied the \emph{existence} of those materials in their possession, and by that denial prevented the tribunal from ever confronting them. \emph{Modnick} itself drew exactly that line (33 Cal.3d at pp.~905--907), holding that concealment-as-denial-of-existence is categorically distinct from perjury-about-a-matter-in-issue and falls on the extrinsic side of \emph{Pico}. The TAC tracks that distinction on its face, which is why a CCP \textsection{} 436 motion cannot reach it: on the well-pleaded facts taken as true, the concealed subject was never ``fairly before the court'' within the meaning of \emph{Modnick}.

\subsection{I.C --- Chronology {\textrightarrow{}} Supreme Court Mapping}\label{i.c-chronology-supreme-court-mapping}

The following table maps each inflection point in the defense's conduct to the Supreme Court holding that treats that conduct as extrinsic fraud:

{\def\LTcaptype{table} % do not increment counter
\begin{longtable}[]{@{}
  >{\raggedright\arraybackslash}p{(\linewidth - 4\tabcolsep) * \real{0.3333}}
  >{\raggedright\arraybackslash}p{(\linewidth - 4\tabcolsep) * \real{0.3333}}
  >{\raggedright\arraybackslash}p{(\linewidth - 4\tabcolsep) * \real{0.3333}}@{}}
\toprule\noalign{}
\begin{minipage}[b]{\linewidth}\raggedright
Date
\end{minipage} & \begin{minipage}[b]{\linewidth}\raggedright
Event
\end{minipage} & \begin{minipage}[b]{\linewidth}\raggedright
Supreme Court Mapping
\end{minipage} \\
\midrule\noalign{}
\endhead
\bottomrule\noalign{}
\endlastfoot
Aug.~18, 2022 & Service of 911 audio by Dolan to defense counsel; Rosario deposition & \emph{Modnick}, 33 Cal.3d at 905 --- concealed ``extrinsic circumstance'' begins \\
May 4, 2023 & Substitution of Attorney CSK {\textrightarrow{}} PSA without conveying production history & \emph{Sanders}, 40 Cal.3d at 615 --- officer-of-court fiduciary breach on file transfer \\
Jan.~17, 2025 & MIL No.~17 filed (``no idea where this 911 call log or audio recording came from'') & \emph{Modnick}, 33 Cal.3d at 907 --- denial of existence \\
Feb.~28, 2025 & Court's institutional-reliance colloquy & \emph{Sanders}, 40 Cal.3d at 615 --- fiduciary reliance extended, then breached \\
Apr.~16, 2025 & Navaratnasingham admits materials ``were provided by Dolan'' (RT 7, p.~213, ll. 13--15) & \emph{Kulchar}, 1 Cal.3d at 472--473 --- diligence discovery trigger \\
Jan.~27 -- Apr.~25, 2025 & Seven consequential rulings (¶ 5 TAC) & \emph{Westphal}, 20 Cal.2d at 397 --- deprivation of opportunity to present \\
\end{longtable}
}

\subsection{I.D --- Downstream Effect on the Court's Rulings}\label{i.d-downstream-effect-on-the-courts-rulings}

Each of the seven consequential rulings was issued inside a record distorted by the concealment, and is therefore a product of the Supreme Court-recognized defect:

\begin{enumerate}
\def\labelenumi{\arabic{enumi}.}
\tightlist
\item
  \textbf{MIL No.~1 Denied} (Jan.~27, p.~41:3--11) --- issued on representations that \emph{Sanders}, 40 Cal.3d at 616, classifies as ``effectively prevent{[}ing{]}'' a fair adversary hearing.
\item
  \textbf{911 CDs Excluded} (Jan.~27, p.~38:22--25) --- same.
\item
  \textbf{911 Transcripts ``Never Admitted''} (Jan.~27, p.~39:18--24) --- categorical exclusion within a \emph{Modnick} 907 record.
\item
  \textbf{Continuance Denied} (Feb.~18, p.~181) --- predicated on ``news to me'' denial now known to be false.
\item
  \textbf{\textsection{} 473 ``Surprise'' Relief Rejected} (Feb.~18, p.~181:20--25) --- same.
\item
  \textbf{Burden-Shift Imposed} (Feb.~28, p.~293:1--3) --- issued during the same colloquy that established the institutional-reliance predicate now breached.
\item
  \textbf{Supplemental Report Excluded} (RT 7, p.~156:24--28) --- issued within the same \emph{Modnick} 907 record.
\end{enumerate}

Under \emph{Sanders}, 40 Cal.3d at 616, and \emph{Modnick}, 33 Cal.3d at 907, these rulings are not ``merits'' adjudications that foreclose equitable vacatur; they are the consequence of the fraud and therefore the very injury the Supreme Court's deprivation rule protects against.

\subsection{I.E --- Step-by-Step Defense Objection Rebuttal (Supreme Court Only)}\label{i.e-step-by-step-defense-objection-rebuttal-supreme-court-only}

{\def\LTcaptype{table} % do not increment counter
\begin{longtable}[]{@{}
  >{\raggedright\arraybackslash}p{(\linewidth - 2\tabcolsep) * \real{0.5000}}
  >{\raggedright\arraybackslash}p{(\linewidth - 2\tabcolsep) * \real{0.5000}}@{}}
\toprule\noalign{}
\begin{minipage}[b]{\linewidth}\raggedright
Defense Objection
\end{minipage} & \begin{minipage}[b]{\linewidth}\raggedright
Supreme Court Foreclosure
\end{minipage} \\
\midrule\noalign{}
\endhead
\bottomrule\noalign{}
\endlastfoot
\emph{Pico}/\emph{Throckmorton} intrinsic ceiling & \emph{Modnick}, 33 Cal.3d at 907 --- denial of existence is extrinsic under California Supreme Court doctrine \\
Lack of diligence & \emph{Kulchar}, 1 Cal.3d at 472--473 --- no duty to discover concealment uniquely within other party's possession until triggering disclosure; \emph{Sanders}, 40 Cal.3d at 615 --- fiduciary breach tolls \\
Pending appeal preclusion & \emph{Mycogen}, 28 Cal.4th at 904; \emph{Boeken}, 48 Cal.4th at 798 --- primary-right analysis turns on injury suffered, not forum; extrinsic-fraud injury is distinct from instructional error \\
Litigation privilege & \emph{Sanders}, 40 Cal.3d at 614 --- equity exception for fraud on the tribunal by officers of the court \\
Non-client standing & \emph{Sanders}, 40 Cal.3d at 615 --- officer-of-court fiduciary duty runs to adjudicative integrity and derivatively to opposing party \\
\emph{Throckmorton} ceiling & \emph{Modnick}, 33 Cal.3d at 905 --- California's broader ``broad concept'' controls; \emph{Throckmorton} is not binding \\
\end{longtable}
}

\subsection{I.F --- The Defense's ``It Got In On Day 7'' Argument, Foreclosed}\label{i.f-the-defenses-it-got-in-on-day-7-argument-foreclosed}

The defense is expected to argue that the ultimate Day-7 admission of the 911 audio cures any prior concealment. That argument is foreclosed by four California Supreme Court holdings. \emph{Westphal}, 20 Cal.2d at 397, fixes the metric at deprivation of ``opportunity to present.'' \emph{Kulchar}, 1 Cal.3d at 471, identifies the injury as denial of ``a fair adversary hearing.'' \emph{Modnick}, 33 Cal.3d at 907, requires that the concealed subject have been ``fairly before the court for adjudication.'' \emph{Sanders}, 40 Cal.3d at 616 (quoting \emph{Larrabee}, 21 Cal.2d at 650--651), measures the harm by conduct that ``effectively prevented \ldots{} a fair adversary hearing.'' Under each of these Supreme Court tests, the inquiry is whether Plaintiff was deprived of the \emph{opportunity to present} his case on a non-corrupted record --- not whether one exhibit was eventually admitted under a corrupted posture. Plaintiff was deprived under each of the four Supreme Court tests. The Day-7 leak-through defense is legally irrelevant to the Supreme Court's deprivation metric.

California's strong policy favoring amendment and resolution on the merits requires that leave be granted. The proposed amendments are made in good faith, cause no prejudice to Defendants, and state viable claims anchored to California Supreme Court authority.

\begin{center}\rule{0.5\linewidth}{0.5pt}\end{center}

\phantomsection\label{toc:ii-statement-of-facts}
\section{II. STATEMENT OF FACTS}\label{ii.-statement-of-facts}

\subsection{A. The Court's Articulated Reliance (February 28, 2025)}\label{a.-the-courts-articulated-reliance-february-28-2025}

On February 28, 2025, the Court explicitly stated the basis for the credibility it had afforded the defense:

\begin{quote}
``I'm not going to question a lawyer's remarks regarding what they have and don't have. They are an officer of the Court, and they are supposed to conduct themselves appropriately before the Court. If they are not, there are other sanctions that they could suffer here. So far I haven't heard or seen anything on the part of the current law firm which tells me that they have conducted themselves improperly.'' --- \emph{The Court (Garrett L. Wong), Hr'g Tr. Feb.~28, 2025, pp.~291:26--292:7.}
\end{quote}

\subsection{B. The Rulings Procured from That Reliance}\label{b.-the-rulings-procured-from-that-reliance}

Based on defense counsel's representations of non-receipt, the Court issued seven consequential rulings:

\begin{enumerate}
\def\labelenumi{\arabic{enumi}.}
\tightlist
\item
  \textbf{MIL No.~1 Denied} (Jan.~27, p.~41:3--11)
\item
  \textbf{911 CDs Excluded} (Jan.~27, p.~38:22--25)
\item
  \textbf{911 Transcripts ``Never Admitted''} (Jan.~27, p.~39:18--24)
\item
  \textbf{Continuance Denied} (Feb.~18, p.~181)
\item
  \textbf{\textsection{} 473 Relief Rejected} (Feb.~18, p.~181:20--25)
\item
  \textbf{Burden-Shift Imposed} (Feb.~28, p.~293:1--3)
\item
  \textbf{Supplemental Report Excluded} (RT 7, p.~156:24--28)
\end{enumerate}

\subsection{C. The Admission That Proved the Misrepresentation (April 16, 2025)}\label{c.-the-admission-that-proved-the-misrepresentation-april-16-2025}

On the seventh day of trial, after months of categorical denials, defense counsel admitted possession:

\begin{quote}
``Sure. They were provided by Dolan. He has had them for whatever, sure.'' --- \emph{Priya D. Navaratnasingham, RT 7 (Apr.~16, 2025), p.~213, ll. 13--15.}
\end{quote}

This admission is irreconcilable with defense counsel's prior representations that they ``never received'' the materials and had ``no idea where'' they came from.

\subsection{D. Procedural History}\label{d.-procedural-history}

\begin{itemize}
\tightlist
\item
  \textbf{December 23, 2025:} Original Complaint filed
\item
  \textbf{February 25, 2026:} First Amended Complaint filed
\item
  \textbf{April 22, 2026:} Defendants filed Motion to Strike under CCP \textsection{} 436
\end{itemize}

\begin{center}\rule{0.5\linewidth}{0.5pt}\end{center}

\phantomsection\label{toc:iii-legal-standard}
\section{III. LEGAL STANDARD}\label{iii.-legal-standard}

\subsection{A. Leave to Amend Should Be Liberally Granted}\label{a.-leave-to-amend-should-be-liberally-granted}

Code of Civil Procedure \textsection{} 473(a)(1) provides that a court ``may, in furtherance of justice, and on any terms as may be proper, allow a party to amend any pleading.''

``The policy favoring amendment is so strong that denial of leave to amend is rarely justified.'' \emph{Howard v. County of San Diego} (2010) 184 Cal.App.4th 1422, 1428.

``If the motion to amend is timely made and the granting of the motion will not prejudice the opposing party, it is error to refuse leave to amend.'' \emph{Morgan v. Superior Court} (1959) 172 Cal.App.2d 527, 530.

\subsection{B. Amendments to Cure Pleading Defects Are Favored}\label{b.-amendments-to-cure-pleading-defects-are-favored}

``Where a demurrer is sustained, leave to amend is routinely granted when there is a reasonable possibility that the pleading can be cured by amendment.'' \emph{Blank v. Kirwan} (1985) 39 Cal.3d 311, 318.

The same principle applies to motions to strike: when a motion to strike is granted based on curable defects, leave to amend should be granted. \emph{Pierson v. Sharp Memorial Hospital, Inc.} (1989) 216 Cal.App.3d 340, 347.

\subsection{C. Prejudice Is the Key Consideration}\label{c.-prejudice-is-the-key-consideration}

``The only legitimate ground for denial of leave to amend is prejudice to the opposing party.'' \emph{Kolani v. Gluska} (1998) 64 Cal.App.4th 402, 412.

Prejudice means actual harm, such as the need to conduct additional discovery or delay in trial, not merely the fact that the opposing party will face an amended pleading. \emph{Atkinson v. Elk Corp.} (2003) 109 Cal.App.4th 739, 761.

\begin{center}\rule{0.5\linewidth}{0.5pt}\end{center}

\phantomsection\label{toc:iv-the-tac-cures-all-eight-defense-objections}
\section{IV. THE TAC CURES ALL EIGHT DEFENSE OBJECTIONS}\label{iv.-the-tac-cures-all-eight-defense-objections}

\subsection{Objection {\textrightarrow{}} Cure Table}\label{objection-cure-table}

{\def\LTcaptype{table} % do not increment counter
\begin{longtable}[]{@{}
  >{\raggedright\arraybackslash}p{(\linewidth - 10\tabcolsep) * \real{0.0323}}
  >{\raggedright\arraybackslash}p{(\linewidth - 10\tabcolsep) * \real{0.2043}}
  >{\raggedright\arraybackslash}p{(\linewidth - 10\tabcolsep) * \real{0.1613}}
  >{\raggedright\arraybackslash}p{(\linewidth - 10\tabcolsep) * \real{0.1075}}
  >{\raggedright\arraybackslash}p{(\linewidth - 10\tabcolsep) * \real{0.1613}}
  >{\raggedright\arraybackslash}p{(\linewidth - 10\tabcolsep) * \real{0.3333}}@{}}
\toprule\noalign{}
\begin{minipage}[b]{\linewidth}\raggedright
\#
\end{minipage} & \begin{minipage}[b]{\linewidth}\raggedright
Defense Objection
\end{minipage} & \begin{minipage}[b]{\linewidth}\raggedright
MPA Page:Line
\end{minipage} & \begin{minipage}[b]{\linewidth}\raggedright
TAC Cure
\end{minipage} & \begin{minipage}[b]{\linewidth}\raggedright
TAC Section/¶
\end{minipage} & \begin{minipage}[b]{\linewidth}\raggedright
Bench Brief/SAC Back-Reference
\end{minipage} \\
\midrule\noalign{}
\endhead
\bottomrule\noalign{}
\endlastfoot
1 & \emph{Kulchar} extrinsic-fraud not viable & --- & Four-channel framework with judicial-reliance quote and seven consequential rulings & \textsection{}I, ¶¶ 1--7 & Bench Brief \textsection{}I.A-I.C \\
2 & Appeal A173827 is remedy & --- & \emph{Kougasian/Stevenot} independent action; \emph{Throckmorton} pre-rebuttal & \textsection{}IV.5, ¶¶ 28--31 & Bench Brief Part VII \\
3 & Civil Code \textsection{} 47(b) bars claims & --- & Nine-Point Evidentiary Chain; \emph{Flatley} illegality; \emph{Action Apartment} extrinsic-fraud exception & \textsection{}II, ¶¶ 8--13 & SAC-06 \textsection{}IV.D.0 \\
4 & Standing / \emph{Shaolian-Norton} & --- & \emph{Estate of Sanders}; \emph{Doctors' Company} aiding-and-abetting framework; \textsection{} 2338 agency & \textsection{}III, ¶¶ 14--21 & Bench Brief \textsection{}II.F \\
5 & Collateral estoppel / res judicata & --- & Five-issues-not-adjudicated block; primary-rights comparison table; \emph{Mycogen/Boeken} test & \textsection{}IV, ¶¶ 22--27 & SAC-06 \textsection{}IV.B \\
6 & CCP \textsection{} 1032 prevailing-party overlap & --- & Express disclaimer; no recovery of underlying costs & \textsection{}V, ¶ 32(b) & --- \\
7 & Damages duplicative of PI case & --- & Reservation paragraph; fraud-delta limitation & \textsection{}V, ¶¶ 32--34 & --- \\
8 & IIED / \emph{Ribas} & --- & No standalone IIED claim; emotional distress as equity factor only & \textsection{}VI, Tier 4, ¶ 38 & --- \\
\end{longtable}
}

\begin{center}\rule{0.5\linewidth}{0.5pt}\end{center}

\phantomsection\label{toc:v-specific-answers-to-defenses-legal-arguments}
\section{V. SPECIFIC ANSWERS TO DEFENSE'S LEGAL ARGUMENTS}\label{v.-specific-answers-to-defenses-legal-arguments}

\subsection{A. Extrinsic Fraud Is Viable Under Kougasian/Olivera}\label{a.-extrinsic-fraud-is-viable-under-kougasianolivera}

Defense argues that Plaintiff cannot maintain an independent action in equity. The TAC cures this by:

\begin{enumerate}
\def\labelenumi{\arabic{enumi}.}
\tightlist
\item
  Pleading the four-channel \emph{Kulchar} framework element-by-element (TAC \textsection{}I, ¶ 3);
\item
  Quoting the Court's February 28, 2025 articulated reliance on officer-of-the-court status (TAC \textsection{}I.5, ¶ 4);
\item
  Documenting the seven consequential rulings procured from that reliance (TAC \textsection{}I.5, ¶ 5);
\item
  Establishing the contradiction between the February 18, 2025 ``news to me'' denial and the April 16, 2025 ``provided by Dolan'' admission (TAC \textsection{}VIII).
\end{enumerate}

Under \emph{Kougasian v. TMSL, Inc.} (2004) 118 Cal.App.4th 1028, 1032--35, an independent action in equity is proper when extrinsic fraud has deprived a party of the opportunity to fully and fairly present his case.

\subsection{B. Appeal Does Not Preclude Independent Action}\label{b.-appeal-does-not-preclude-independent-action}

Defense argues that Plaintiff's remedy is the pending appeal (A173827). The TAC cures this by:

\begin{enumerate}
\def\labelenumi{\arabic{enumi}.}
\tightlist
\item
  Distinguishing the appeal (instructional error) from this action (institutional fraud) (TAC \textsection{}IV.5, ¶¶ 28--29);
\item
  Citing \emph{Kougasian} and \emph{In re Marriage of Stevenot} (1984) 154 Cal.App.3d 1051 for the principle that an independent equitable action is cognizable alongside a direct appeal;
\item
  Pre-rebutting defense reliance on \emph{Throckmorton} as non-binding federal authority superseded by California's broader doctrine (TAC \textsection{}II, ¶ 13).
\end{enumerate}

\subsection{C. Litigation Privilege Does Not Apply}\label{c.-litigation-privilege-does-not-apply}

Defense invokes Civil Code \textsection{} 47(b). The TAC cures this by:

\begin{enumerate}
\def\labelenumi{\arabic{enumi}.}
\tightlist
\item
  Documenting the Nine-Point Evidentiary Chain showing 893 days of defense possession (TAC \textsection{}II, ¶ 9);
\item
  Pleading non-communicative conduct (silence, file-transfer concealment, failure to correct) that falls outside the privilege (TAC \textsection{}II, ¶ 8);
\item
  Invoking the \emph{Flatley v. Mauro} (2006) 39 Cal.4th 299 illegality exception (Bus. \& Prof.~Code \textsection{} 6128(a) as criminal conduct) (TAC \textsection{}II, ¶¶ 10--11);
\item
  Citing \emph{Action Apartment Ass'n v. City of Santa Monica} (2007) 41 Cal.4th 1232 for the extrinsic-fraud carveout (TAC \textsection{}II, ¶ 12).
\end{enumerate}

\subsection{D. Standing Is Established}\label{d.-standing-is-established}

Defense cites \emph{Shaolian} and \emph{Norton} for lack of standing. The TAC cures this by:

\begin{enumerate}
\def\labelenumi{\arabic{enumi}.}
\tightlist
\item
  Distinguishing \emph{Shaolian} (direct actions for policy benefits) from equitable fraud-on-the-court theories (TAC \textsection{}III, ¶ 15);
\item
  Citing \emph{Estate of Sanders} (1985) 40 Cal.3d 607 (attorneys who participate in fraud on the court are proper parties) (TAC \textsection{}III, ¶ 15(c));
\item
  Pleading agency/ratification under Civil Code \textsection{} 2338 (TAC \textsection{}III, ¶¶ 16--17);
\item
  Invoking \emph{Doctors' Company v. Superior Court} (1989) 49 Cal.3d 39 for insurer aiding-and-abetting liability (TAC \textsection{}III, ¶ 18).
\end{enumerate}

\subsection{E. Collateral Estoppel Does Not Apply}\label{e.-collateral-estoppel-does-not-apply}

Defense argues res judicata/collateral estoppel. The TAC cures this by:

\begin{enumerate}
\def\labelenumi{\arabic{enumi}.}
\tightlist
\item
  Enumerating five issues never adjudicated in the underlying trial (TAC \textsection{}IV, ¶ 24);
\item
  Providing a primary-rights comparison table under \emph{Mycogen} and \emph{Boeken} (TAC \textsection{}IV, ¶ 25);
\item
  Applying the \emph{Lucido v. Superior Court} (1990) 51 Cal.3d 335 factors (TAC \textsection{}IV, ¶ 27).
\end{enumerate}

\subsection{F. Damages Are Properly Limited}\label{f.-damages-are-properly-limited}

Defense argues damages are duplicative. The TAC cures this by:

\begin{enumerate}
\def\labelenumi{\arabic{enumi}.}
\tightlist
\item
  Including express disclaimers of PI damages, prevailing-party costs, duplicative 802 damages, and pro-per fees (TAC \textsection{}V, ¶ 32);
\item
  Limiting compensatory relief to ``fraud-delta'' costs---incremental expenses caused by the fraud (TAC \textsection{}VI, Tier 2, ¶ 36).
\end{enumerate}

\subsection{G. Extrinsic Fraud Measured by Deprivation, Not Admission}\label{g.-extrinsic-fraud-measured-by-deprivation-not-admission}

The defense is expected to argue that the Day-7 admission of the 911 audio cures the concealment. That argument is foreclosed by four California Supreme Court holdings that independently locate the operative injury at the deprivation of opportunity, not at the admission of evidence:

\begin{enumerate}
\def\labelenumi{\arabic{enumi}.}
\tightlist
\item
  \emph{Westphal v. Westphal} (1942) 20 Cal.2d 393, 397 --- extrinsic fraud exists where a party is ``prevented from exhibiting fully his case.''
\item
  \emph{Kulchar v. Kulchar} (1969) 1 Cal.3d 467, 471 --- the ``primary right'' protected is ``a fair adversary hearing.''
\item
  \emph{In re Marriage of Modnick} (1983) 33 Cal.3d 897, 907 --- the concealed subject must have been ``fairly before the court for adjudication''; a late-admitted exhibit does not retroactively cure the rulings that issued within the corrupted record.
\item
  \emph{Estate of Sanders} (1985) 40 Cal.3d 607, 616 (quoting \emph{Larrabee v. Tracy} (1943) 21 Cal.2d 645, 650--651) --- the measure is whether the conduct ``effectively prevented \ldots{} a fair adversary hearing.''
\end{enumerate}

Under each test, Plaintiff was deprived of the opportunity to present his case on a non-corrupted record. The Court's seven consequential rulings issued inside that corrupted record (Memo \textsection{} I.D). The TAC pleads that deprivation in \textsection{}IV.5.5 and ¶¶ 34.5--34.6, mirroring these four Supreme Court holdings. The ``it got in on Day 7'' defense is legally irrelevant to the Supreme Court's deprivation metric.

\begin{center}\rule{0.5\linewidth}{0.5pt}\end{center}

\phantomsection\label{toc:vi-no-prejudice-to-defendants}
\section{VI. NO PREJUDICE TO DEFENDANTS}\label{vi.-no-prejudice-to-defendants}

The proposed amendments do not prejudice Defendants because:

\begin{enumerate}
\def\labelenumi{\arabic{enumi}.}
\item
  \textbf{No New Parties:} The TAC names the same defendants as the operative complaint (with corrected identification of the defense firm and individual attorneys).
\item
  \textbf{No New Theories:} The TAC pleads the same claims---extrinsic fraud, constructive fraud, \textsection{} 6128(a) violations---with greater specificity.
\item
  \textbf{No Additional Discovery Needed:} The TAC is based on the certified record already in Defendants' possession.
\item
  \textbf{No Trial Delay:} This motion is made early in the proceedings.
\end{enumerate}

\begin{center}\rule{0.5\linewidth}{0.5pt}\end{center}

\phantomsection\label{toc:vii-conclusion}
\section{VII. CONCLUSION}\label{vii.-conclusion}

For the foregoing reasons, Plaintiff respectfully requests that this Court grant leave to file the Third Amended Complaint in the form lodged herewith.

\begin{center}\rule{0.5\linewidth}{0.5pt}\end{center}

\textbf{Dated:} \_\_\_\_\_\_\_\_\_\_\_\_\_\_\_

\textbf{FRANCISCUS DYLAN ROSARIO}\\
Plaintiff, In Pro Per\\
7624 Melody Drive\\
Rohnert Park, CA 94928\\
Tel: 650.488.7510\\
Email: soltrinox@gmail.com

\nolinenumbers
\input{TAC-801-SIGNATURE-BLOCK.tex}

\end{document}
